\section*{Hoofdstuk \ref{ch:g6:latitude-longitude} --- Breedte, lengte en tijdzones}

\begin{solution}{exo:g6:latitude-longitude:1}
``De \omterm{def:g6:latitude-longitude:latitude}{breedtegraad} wordt gemeten vanaf de \omterm{def:g3:continents-and-oceans:equator}{evenaar} en loopt van $0^\circ$ tot
$90^\circ$ (noord of zuid); de \omterm{def:g6:latitude-longitude:longitude}{lengtegraad} wordt gemeten vanaf de
\omterm{def:g6:latitude-longitude:longitude}{nulmeridiaan} en loopt van $0^\circ$ tot $180^\circ$ (oost of west).''
\end{solution}

\begin{solution}{exo:g6:latitude-longitude:2}
De \omterm{def:g3:continents-and-oceans:poles}{Noordpool}: $90^\circ$\,NB. De \omterm{def:g3:continents-and-oceans:poles}{Zuidpool}: $90^\circ$\,ZB. Een stad op de
\omterm{def:g3:continents-and-oceans:equator}{evenaar}: $0^\circ$.
\end{solution}

\begin{solution}{exo:g6:latitude-longitude:3}
Sydney: \omterm{def:g3:continents-and-oceans:equator}{zuidelijk halfrond}, oost. New York: \omterm{def:g3:continents-and-oceans:equator}{noordelijk halfrond}, west.
Caïro: \omterm{def:g3:continents-and-oceans:equator}{noordelijk halfrond}, oost. Rio de Janeiro: \omterm{def:g3:continents-and-oceans:equator}{zuidelijk halfrond},
west.
\end{solution}

\begin{solution}{exo:g6:latitude-longitude:4}
(a) Waar. (b) Niet waar: de \omterm{def:g6:latitude-longitude:latitude}{breedtecirkels} worden korter naar de \omterm{def:g3:continents-and-oceans:poles}{polen}
toe; de langste is de \omterm{def:g3:continents-and-oceans:equator}{evenaar}. (c) Waar. (d) Niet waar: ze loopt door
Greenwich, in Londen.
\end{solution}

\begin{solution}{exo:g6:latitude-longitude:5}
Caïro. New York. Tokio ligt op ongeveer $36^\circ$\,NB, $140^\circ$\,OL
(iets boven de \omterm{def:g6:latitude-longitude:latitude}{breedtecirkel} $30^\circ$\,NB, iets voorbij de \omterm{def:g6:latitude-longitude:longitude}{meridiaan}
$120^\circ$\,OL, nog voor $150^\circ$\,OL).
\end{solution}

\begin{solution}{exo:g6:latitude-longitude:6}
$45 \div 15 = 3$: hun middagen liggen 3~uur uit elkaar. Het eerst is het
middag in stad~B, die oostelijker ligt.
\end{solution}

\begin{solution}{exo:g6:latitude-longitude:7}
$73 \div 15 \approx 4.9$, afgerond 5: de zone is $+5$. Als het in
Greenwich twaalf uur 's middags is, is het in Mumbai $12 + 5 = 17$, dus
17~uur.
\end{solution}

\begin{solution}{exo:g6:latitude-longitude:8}
$30 \times 111 = 3\,330$~km. De \omterm{def:g3:continents-and-oceans:poles}{Noordpool} ligt op $90^\circ$:
$90 \times 111 = 9\,990$~km, ongeveer 10\,000~km.
\end{solution}

\begin{solution}{exo:g6:latitude-longitude:9}
$9 - (-5) = 14$~uur verschil. $8 + 14 = 22$: het is maandag 22~uur in
Tokio.
\end{solution}

\begin{solution}{exo:g6:latitude-longitude:10}
Ze heeft gelijk. De Aarde draait van west naar oost, dus de steden in het
oosten worden het eerst naar de Zon toe gedraaid: ze zien haar opkomen, op
haar hoogst staan en ondergaan voor de steden verder naar het westen.
\end{solution}

\begin{solution}{exo:g6:latitude-longitude:11}
(a) $42 \times 3 = 126$~seconden, iets meer dan 2~minuten. (b)
$126 \div 240 \approx 0.5$: ongeveer een halve \omterm{def:g6:latitude-longitude:longitude}{lengtegraad}. (c)
$0.5 \times 111 \approx 55$~km (met $0.525$: ongeveer 58~km). Een fout van
ongeveer vijftig kilometer: genoeg om een eiland te missen, of op een \omterm{def:g1:landscapes-around-us:water}{kust}
te lopen waarvan je dacht dat die nog ver weg was.
\end{solution}

\begin{solution}{exo:g6:latitude-longitude:12}
(a) $18 - 9 = 9$: maandag 9~uur 's ochtends in Greenwich. (b)
$9 + 10 = 19$: maandag 19~uur in Greenwich. (c) $19 - 8 = 11$: maandag
11~uur 's ochtends in Los Angeles. Het vliegtuig vertrok op maandag om
18~uur (tijd in Tokio) en landt om 11~uur op dezelfde maandag (tijd in Los
Angeles): de klokken bij aankomst wijzen een vroeger uur aan dan de
klokken bij vertrek, omdat het vliegtuig de \omterm{def:g6:latitude-longitude:date-line}{datumgrens} naar het oosten is
overgestoken en een dag is teruggegaan.
\end{solution}

\begin{solution}{pb:g6:latitude-longitude:1}
\textbf{1.} Ana (Rio de Janeiro, $23^\circ$\,ZB) en Mia (Auckland,
$37^\circ$\,ZB).

\textbf{2.} Ana en Sam ($43^\circ$\,WL en $118^\circ$\,WL). Lena woont op de
\omterm{def:g6:latitude-longitude:longitude}{nulmeridiaan}.

\textbf{3.} Lena: Londen ligt op $52^\circ$\,NB, verder van de \omterm{def:g3:continents-and-oceans:equator}{evenaar} dan
alle anderen.

\textbf{4.} Mia: Auckland ligt op $175^\circ$\,OL, maar $180 - 175 = 5$
graden van de \omterm{def:g6:latitude-longitude:longitude}{meridiaan} van $180^\circ$.

\textbf{5.} $360^\circ \div 24 = 15^\circ$.

\textbf{6.} Londen: $0$. Rio de Janeiro: $43 \div 15 \approx 2.9$, west,
dus $-3$. Tokio: $140 \div 15 \approx 9.3$, oost, dus $+9$. Auckland:
$175 \div 15 \approx 11.7$, oost, dus $+12$. Los Angeles:
$118 \div 15 \approx 7.9$, west, dus $-8$.

\textbf{7.} Rio de Janeiro: $20 - 3 = 17$, vrijdag 17~uur. Tokio:
$20 + 9 = 29$; $29 - 24 = 5$: zaterdag 5~uur 's ochtends.

\textbf{8.} Tokio ligt ver naar het oosten: zijn klokken lopen 9~uur voor op
die van Londen. Om 20~uur in Londen is het in Tokio al na middernacht, dus
is zijn zaterdag begonnen.

\textbf{9.} $20 - 8 = 12$: vrijdag twaalf uur 's middags.

\textbf{10.} $12 - (-8) = 20$~uur voor.

\textbf{11.} Als het schip de \omterm{def:g6:latitude-longitude:date-line}{datumgrens} naar het oosten oversteekt, gaat
het terug naar de dag ervoor.

\textbf{12.} $20 + 12 = 32$; $32 - 24 = 8$: als het gesprek begint, is het
in Auckland \textbf{zaterdag 8~uur 's ochtends}. Mia doet mee aan het
gesprek bij het ontbijt, terwijl Sam op vrijdag aan de lunch zit.
\end{solution}
